\chapter{Simulation Pipeline}\label{app:pipeline}

\section{Pipeline Architecture}

Figure~\ref{fig:architecture} shows the full pipeline architecture, from local configuration to HPC output.

\begin{figure}[ht]
  \centering
  \includegraphics[width=\textwidth]{img/architecture.png}
  \caption{Automated simulation pipeline architecture. Local machine (blue) deploys to the ETH Euler HPC cluster. The login node builds all FE models; compute nodes run the solver and post-processing jobs in parallel under SLURM. The FLC aggregation job runs after all solver jobs complete via a SLURM \texttt{afterok} dependency.}
  \label{fig:architecture}
\end{figure}

\section{Repository Structure}

The project is hosted on GitHub (\texttt{andrecruzf/AbaqusProject}). Key files:

\begin{tabbing}
 \hspace*{4.5cm} \= \kill
 \texttt{config.py}             \> Central configuration — all model parameters \\[0.5ex]
 \texttt{build\_model.py}       \> Pipeline orchestrator, calls all modules sequentially \\[0.5ex]
 \texttt{modules/}              \> \texttt{parts}, \texttt{assembly}, \texttt{material}, \texttt{step}, \texttt{contact}, \texttt{boundary}, \texttt{output}, \texttt{job} \\[0.5ex]
 \texttt{VUMAT\_explicit.f}     \> Hosford-Coulomb VUMAT subroutine (Fortran 77) \\[0.5ex]
 \texttt{geometries/}           \> Reference blank geometry files (W20--W200) \\[0.5ex]
 \texttt{postproc.py}           \> Extracts strain path and energy ratio from ODB \\[0.5ex]
 \texttt{postproc\_movie.py}    \> Renders SDV1 (EQPS) animation to \texttt{.webm} \\[0.5ex]
 \texttt{flc\_plot.py}          \> Aggregates all \texttt{strain\_path.csv} files and plots FLC \\[0.5ex]
 \texttt{deploy\_all.sh}        \> Entry point: push scripts to Euler, build all widths, submit jobs \\[0.5ex]
 \texttt{run\_cluster.sh}       \> SLURM job script: solver + postproc + copy results to home \\[0.5ex]
 \texttt{run\_flc.sh}           \> SLURM FLC aggregation job (runs after all solver jobs)
\end{tabbing}

\section{Configurable Simulation Parameters}
\label{app:configuration_parameters}

The numerical framework is controlled from a central configuration file. The
main text describes the modelling assumptions and data flow, while the
configuration groups are collected here to keep the numerical model chapter
focused on the general framework.

\begin{table}[ht]
\begin{center}
\caption{Main user-tunable parameter groups in the numerical framework.}
\label{tab:framework_parameters}
\begin{tabular}{@{}lll@{}}\toprule[1.5pt]
Parameter group & Examples & Purpose \\ \midrule
Test configuration & Nakazima, Marciniak, Punch-in-Punch & Select forming setup \\
Specimen geometry & Width, thickness, geometry source & Build the tested blank \\
Tool geometry & Punch radius, punch variant, die radii & Study tooling and size effects \\
Process conditions & Punch travel, speed, friction & Match or vary test conditions \\
Material orientation & Rolling/transverse angle & Evaluate anisotropy effects \\
Mesh control & Refinement factor, seed mode, thickness seeds & Control spatial resolution \\
Explicit settings & Mass scaling, output frequency, CPUs & Balance cost and accuracy \\
Post-processing & Dome radius, fracture detector, V\&H thresholds & Standardise limit extraction \\ \bottomrule[1.5pt]
\end{tabular}
\end{center}
\end{table}

For the mass-scaling and mesh-refinement sensitivity study, the framework
generates a fully factorial grid of simulations while keeping the model
definition and post-processing schema unchanged. The grid used in the current
study is reported in Table~\ref{tab:study_grid}.

\begin{table}[ht]
\begin{center}
\caption{Parameter grid for the mass scaling $\times$ mesh refinement sensitivity study (Nakazima W200, $t=1.75$\,mm). Each cell is one independent simulation.}
\label{tab:study_grid}
\begin{tabular}{@{}lccc@{}}\toprule[1.5pt]
 & \multicolumn{3}{c}{Mesh refinement factor $f_\mathrm{MR}$} \\
\cmidrule(l){2-4}
$\Delta t$ [s] & 1 & 2 & 4 \\ \midrule
$1\times10^{-7}$ & \checkmark & \checkmark & \checkmark \\
$1\times10^{-6}$ & \checkmark & \checkmark & \checkmark \\
$1\times10^{-5}$ & \checkmark & \checkmark & \checkmark \\
$1\times10^{-4}$ & \checkmark & \checkmark & \checkmark \\
\bottomrule[1.5pt]
\end{tabular}
\end{center}
\end{table}

\section{Running a Simulation Campaign}

To launch a full campaign for all widths with a given test configuration:
\lstset{language=bash,numbers=none}
\begin{lstlisting}[frame=lines]
# Deploy and submit all widths (defaults to TEST_TYPE and thickness from config.py):
./deploy_all.sh

# Or specify test type, thickness, and orientation angle explicitly:
./deploy_all.sh nakazima 1.0 0
./deploy_all.sh pip 1.0 0
\end{lstlisting}

Results are copied from scratch to the home directory after each job, and the FLC aggregation job runs automatically after all solver jobs complete.

\section{Cluster Configuration}

Simulations run on the ETH Euler cluster (\texttt{acruzfaria@euler.ethz.ch}) using SLURM. Each solver job requests 24 CPUs, 64\,GB RAM, and a 24-hour time limit. The Abaqus module is loaded via \texttt{module load abaqus/2023}. Post-processing scripts use \texttt{xvfb-run} for headless ODB rendering. Results are stored in \texttt{/cluster/scratch/acruzfaria/} during the run and copied to \texttt{/cluster/home/acruzfaria/} on completion.
